%------------------------------------------------------------------------------%
\begin{question}
  Determine os autovetores da matriz
  \(
    A =
    \begin{bmatrix}
      4 & 1 \\
      2 & 3
    \end{bmatrix}
  \)

  \bigskip
  Note que, no exercício anterior, determinamos que os autovalores de $A$ são
  \[
    \lambda_1 = 5
    \qquad
    \lambda_2 = 2
  \]
\end{question}

%------------------------------------------------------------------------------%
\begin{resolution}{Autovetores Associados a \(\lambda_1 = 5\)}
\begin{columns}
\column{0.5\textwidth}
  \begin{align*}
    \onslide<+->{(A-\lambda_1 I)\vec{v} & = \vec{0} \\}
    \onslide<+->{(A-5I)\vec{v} & = \vec{0}}
    \\
    \onslide<+->{\left(
      \begin{bmatrix}
        4 & 1 \\
        2 & 3
      \end{bmatrix}
      -
      \begin{bmatrix}
        5 & 0 \\
        0 & 5
      \end{bmatrix}
    \right)
        \begin{bmatrix} x \\ y \end{bmatrix}
    & = \begin{bmatrix} 0 \\ 0 \end{bmatrix}}
    \\
    \onslide<+->{\begin{bmatrix}
      -1  &  1  \\
       2  & -2
    \end{bmatrix}
        \begin{bmatrix} x \\ y \end{bmatrix}
    & = \begin{bmatrix} 0 \\ 0 \end{bmatrix}}
    \\
    \onslide<+->{\begin{cases}
      -x +  y = 0 \\
      2x - 2y = 0
    \end{cases}}
  \end{align*}
\column{0.5\textwidth}
  \begin{align*}
    \onslide<+->{-x + y & = 0 \\}
    \onslide<+->{     y & = x}
  \end{align*}
  \onslide<+->{Assim}
  \(
    \onslide<+->{\vec{v}}
    \onslide<+->{=   \begin{bmatrix} x \\ x \end{bmatrix}}
    \onslide<+->{= x \begin{bmatrix} 1 \\ 1 \end{bmatrix}}
  \)

  \medskip
  \onslide<+->{Os autovetores são \\
  os múltiplos não nulos}
  \onslide<+->{de
  \[
    \vec{v}_1 =
    \begin{bmatrix}
      1 \\
      1
    \end{bmatrix}
  \]}
\end{columns}
\end{resolution}


%------------------------------------------------------------------------------%
\begin{resolution}{Autovetores Associados a \(\lambda_2 = 2\)}
\begin{columns}
\column{0.5\textwidth}
  \begin{align*}
    \onslide<+->{(A-\lambda_2 I)\vec{v} & = \vec{0} \\}
    \onslide<+->{(A-2I)\vec{v} & = \vec{0} }
    \\
    \onslide<+->{\left(
      \begin{bmatrix}
        4 & 1 \\
        2 & 3
      \end{bmatrix}
      -
      \begin{bmatrix}
        2 & 0 \\
        0 & 2
      \end{bmatrix}
    \right)
        \begin{bmatrix} x \\ y \end{bmatrix}
    & = \begin{bmatrix} 0 \\ 0 \end{bmatrix}}
    \\
    \onslide<+->{\begin{bmatrix}
      2 & 1  \\
      2 & 1
    \end{bmatrix}
        \begin{bmatrix} x \\ y \end{bmatrix}
    & = \begin{bmatrix} 0 \\ 0 \end{bmatrix}}
    \\
    \onslide<+->{\begin{cases}
      2x + y = 0 \\
      2x + y = 0
    \end{cases}}
  \end{align*}
\column{0.5\textwidth}
  \[
    \onslide<+->{y = -2x}
  \]
  \onslide<+->{Assim}
  \(
    \onslide<+->{\vec{v}}
    \onslide<+->{= \begin{bmatrix}
      x \\
      -2x
    \end{bmatrix}}
    \onslide<+->{= x
    \begin{bmatrix}
      1 \\
      -2
    \end{bmatrix}}
  \)

  \medskip
  \onslide<+->{Os autovetores são \\
  os múltiplos não nulos}
  \onslide<+->{de
  \[
    \vec{v}_2 =
    \begin{bmatrix}
      1 \\
      -2
    \end{bmatrix}
  \]}
\end{columns}
\end{resolution}

%------------------------------------------------------------------------------%
\endinput


